\documentclass[a4paper,11pt]{article}
\usepackage[utf8]{inputenc}
\usepackage{geometry}
\usepackage{xcolor}
\usepackage{listings}
\usepackage{hyperref}
\usepackage{graphicx}
\usepackage{amsmath}
\usepackage{amssymb}
\usepackage{tcolorbox}
\usepackage{float}

% Geometria della pagina
\geometry{top=2.5cm, bottom=2.5cm, left=2.5cm, right=2.5cm}

% Colori custom
\definecolor{ibmblue}{RGB}{0, 45, 156}
\definecolor{ibmred}{RGB}{250, 75, 76}
\definecolor{codegreen}{rgb}{0,0.6,0}
\definecolor{codegray}{rgb}{0.5,0.5,0.5}
\definecolor{codepurple}{rgb}{0.58,0,0.82}
\definecolor{backcolour}{rgb}{0.95,0.95,0.92}

% Setup per il codice Python
\lstdefinestyle{mystyle}{
    backgroundcolor=\color{backcolour},   
    commentstyle=\color{codegreen},
    keywordstyle=\color{magenta},
    numberstyle=\tiny\color{codegray},
    stringstyle=\color{codepurple},
    basicstyle=\ttfamily\footnotesize,
    breakatwhitespace=false,         
    breaklines=true,                 
    captionpos=b,                    
    keepspaces=true,                 
    numbers=left,                    
    numbersep=5pt,                  
    showspaces=false,                
    showstringspaces=false,
    showtabs=false,                  
    tabsize=2
}
\lstset{style=mystyle}

% Box personalizzati
\newtcolorbox{studentbox}{colback=blue!5!white,colframe=blue!75!black,title=Student Activity}
\newtcolorbox{theorybox}{colback=gray!5!white,colframe=gray!50!black,title=Theory Recap}
\newtcolorbox{warningbox}{colback=yellow!10!white,colframe=orange!80!black,title=Important: The Black Box (Oracle)}
\newtcolorbox{composerbox}{colback=purple!5!white,colframe=purple!75!black,title=IBM Composer View (Q-Sphere)}
\newtcolorbox{deepdivebox}{colback=teal!5!white,colframe=teal!75!black,title=Deep Dive: Beyond Deutsch (QC14)}
\newtcolorbox{mathbox}{colback=red!5!white,colframe=red!75!black,title=Step-by-Step State Evolution}

\title{\textbf{Laboratory Guide 3: Quantum Algorithms}\\
\large Oracles, Phase Kickback, and the Quantum Advantage}
\author{Course: Quantum Computing Foundation}
\date{}

\begin{document}

\maketitle

\section{Introduction and Objectives}
In the previous labs, we learned how to manipulate qubits and create entanglement. Now, we transition from physics to computer science. We will implement \textbf{Deutsch's Algorithm}, a milestone in quantum computing. It is the simplest algorithm that unequivocally proves a quantum computer can solve a specific problem faster than any classical machine.

\vspace{0.3cm}
\noindent \textbf{Prerequisites:}
\begin{itemize}
    \item Lesson QC12: \textit{Algorithms - Quantum Parallelism}
    \item Lesson QC13: \textit{Deutsch's Algorithm}
    \item Lab 2: \textit{Multi-Qubit Systems}
\end{itemize}

\noindent \textbf{Learning Objectives:}
\begin{enumerate}
    \item Understand the Query Model and the role of a Quantum Oracle.
    \item Visualize and master the "Phase Kickback" mechanism using the Composer.
    \item Track the mathematical evolution of the quantum state step-by-step.
    \item Implement Deutsch's Algorithm in Qiskit using professional coding standards.
\end{enumerate}

\begin{warningbox}
\textbf{What is an Oracle?}\\
In computational complexity, an Oracle is a "Black Box". Suppose we have a function $f(x)$. We can give the Oracle an input $x$, and it returns the output $f(x)$. Crucially, we \textbf{cannot} look inside the box to see the circuitry or the code. Our goal is to discover global properties of $f(x)$ while minimizing the number of queries (questions) asked to the Oracle.
\end{warningbox}

\newpage

\section{Part A: The Secret Mechanism (IBM Quantum Composer)}
\textit{Platform: IBM Quantum Composer}

\subsection{Exercise 1: Phase Kickback}
\textbf{Context:} As explained in Lesson \textbf{QC13}, quantum parallelism allows us to evaluate $f(0)$ and $f(1)$ simultaneously. However, measuring the state directly destroys the superposition and only gives us one answer. To extract the relation between $f(0)$ and $f(1)$, we rely on a physical phenomenon called \textbf{Phase Kickback}. 

When a target qubit is in the superposition state $|-\rangle$, applying a controlled operation (like a CNOT) does not flip the target's probabilities. Instead, the phase change "kicks back" and alters the phase of the \textit{control} qubit.

\begin{studentbox}
\textbf{Action:}
\begin{enumerate}
    \item Add an \texttt{H} gate to $q_0$ (Control). It is now in state $|+\rangle$.
    \item Add an \texttt{X} gate, then an \texttt{H} gate to $q_1$ (Target). It is now in state $|-\rangle$.
    \item Apply a \texttt{CNOT} gate (Control $q_0$, Target $q_1$).
    \item \textbf{Look at the Q-Sphere} before and after the CNOT.
\end{enumerate}
\textbf{Expected Effect:}
Normally, a CNOT flips the target. Here, the target $q_1$ remains strictly $|-\rangle$. Instead, the phase of the control qubit $q_0$ flips from $|+\rangle$ to $|-\rangle$! The "action" was transferred backward to the control.
\end{studentbox}

\begin{composerbox}
\textbf{Q-Sphere Visualization}\\
Before the CNOT, the dots on the Q-Sphere for the control qubit are Blue (Positive phase). After the CNOT, the dot corresponding to the $|1\rangle$ component of the control qubit turns \textbf{Red} (Negative phase). The target qubit's probabilities and phases did not change at all. This Red dot is the core of quantum computation.
\end{composerbox}

\newpage

\section{Part B: The Algorithm Step-by-Step}

\subsection{Exercise 2: Finding the Function Type}
\textbf{The Problem:} We are given a function $f(x)$ that takes 1 bit (0 or 1) as input and returns 1 bit. 
There are two possible categories for this function:
\begin{itemize}
    \item \textbf{Constant:} Returns the same value for all inputs (i.e., $f(0) = f(1)$).
    \item \textbf{Balanced:} Returns different values (i.e., $f(0) \neq f(1)$).
\end{itemize}
Classically, you must evaluate the function twice ($x=0$ and $x=1$) to compare the outputs. Deutsch's Algorithm solves this with \textbf{exactly 1 quantum query}.

\begin{mathbox}
\textbf{How it works: The Mathematical Evolution}\\
Let's trace the state $|\psi\rangle$ of our two qubits (Input $x$, Ancilla $y$) through the algorithm.
\begin{enumerate}
    \item \textbf{Initialization:} We start with the input at $|0\rangle$ and the ancilla at $|1\rangle$.
    $$|\psi_0\rangle = |0\rangle \otimes |1\rangle$$
    
    \item \textbf{Superposition (Apply H to both):} We prepare the input to query both values simultaneously, and the ancilla to absorb the phase kickback.
    $$|\psi_1\rangle = |+\rangle \otimes |-\rangle = \left( \frac{|0\rangle + |1\rangle}{\sqrt{2}} \right) \left( \frac{|0\rangle - |1\rangle}{\sqrt{2}} \right)$$
    
    \item \textbf{The Oracle Query $U_f$:} The quantum oracle performs the operation $|x\rangle|y\rangle \to |x\rangle|y \oplus f(x)\rangle$. Because the ancilla is in the $|-\rangle$ state, the function's output $f(x)$ is transferred to the phase of the input qubit:
    $$|\psi_2\rangle = \left( \frac{(-1)^{f(0)}|0\rangle + (-1)^{f(1)}|1\rangle}{\sqrt{2}} \right) |-\rangle$$
    
    \item \textbf{Interference (Apply H to Input):} We apply a final Hadamard gate to the input qubit to extract the phase difference.
    \begin{itemize}
        \item If $f(0) = f(1)$ (Constant), the phases are the same, and the state interferes constructively back to $\pm|0\rangle$.
        \item If $f(0) \neq f(1)$ (Balanced), the phases are opposite, and the state interferes constructively to $\pm|1\rangle$.
    \end{itemize}
\end{enumerate}
\textbf{Conclusion:} A measurement of $0$ means Constant. A measurement of $1$ means Balanced.
\end{mathbox}

\newpage

\subsection{Exercise 3: Coding Deutsch's Algorithm in Qiskit}
We will now implement the algorithm described above. We will use a CNOT gate as our Oracle. A CNOT acts as $f(x) = x$, which is a \textbf{Balanced} function.

\begin{tcolorbox}[colback=green!5!white,colframe=green!75!black,title=Coding Task: Deutsch's Algorithm Implementation]
\begin{lstlisting}[language=Python]
"""
Deutsch's Algorithm Implementation.
This script demonstrates the algorithm using a balanced oracle.
"""
from qiskit import QuantumCircuit
from qiskit_aer import AerSimulator
from qiskit.visualization import plot_histogram

def run_deutsch_algorithm() -> dict:
    """
    Builds and executes Deutsch's algorithm for a balanced function f(x) = x.
    
    Returns:
        dict: The measurement counts, proving the function type.
    """
    # 1. Setup: 1 Input Qubit (q0), 1 Ancilla Qubit (q1), 1 Classical Bit
    qc = QuantumCircuit(2, 1)

    # 2. State Preparation (Superposition and Phase Target)
    qc.h(0)           # Input to |+>
    qc.x(1)           # Ancilla to |1>
    qc.h(1)           # Ancilla to |-> for phase kickback
    qc.barrier()

    # 3. The Oracle: Balanced function (CNOT)
    # This represents the unknown black box.
    qc.cx(0, 1)
    qc.barrier()

    # 4. Interference and Measurement
    qc.h(0)           # Extract the relative phase
    qc.measure(0, 0)  # Read the result into the classical bit

    # 5. Execution
    simulator = AerSimulator()
    job = simulator.run(qc, shots=1000)
    return job.result().get_counts()

# Run the simulation and print the results
counts = run_deutsch_algorithm()
print(f"Algorithm Output: {counts}")
\end{lstlisting}
\textbf{Analysis:} Because the oracle we built is balanced, the output will be exactly `{'1': 1000}`. If you replaced the Oracle block with an empty space (Identity gate), it would act as $f(x) = 0$ (Constant), and the output would be `{'0': 1000}`.
\end{tcolorbox}

\begin{deepdivebox}
\textbf{Beyond Deutsch: Why does this matter? (Lesson QC14)}\\
Deutsch's algorithm has no direct commercial application. However, it proves that quantum interference can solve a problem in $\mathcal{O}(1)$ steps instead of $\mathcal{O}(N)$.
\begin{itemize}
    \item \textbf{Deutsch-Jozsa:} Expands this to $N$ bits. A classical computer needs $2^{N-1} + 1$ queries in the worst case. The quantum computer still needs exactly \textbf{1 query}. This represents an exponential speedup.
    \item \textbf{Grover \& Shor:} The logic of Oracles and Phase Kickback is the exact same engine used by Grover's algorithm to search databases and Shor's algorithm to break RSA encryption.
\end{itemize}
\end{deepdivebox}

\newpage

\section{Part C: Assessment (Quiz)}

\textbf{Q1. The Power of the Oracle}\\
In the Query Model of computation (QC12), what is the primary metric used to evaluate the efficiency of a quantum algorithm?
\begin{itemize}
    \item[A)] The clock speed of the quantum processor.
    \item[B)] The physical number of qubits required.
    \item[C)] The number of times the algorithm evaluates the Oracle function $f(x)$.
\end{itemize}

\vspace{0.3cm}

\textbf{Q2. The Core Mechanism}\\
In the mathematical evolution of Deutsch's algorithm, what causes the value of $f(x)$ to be encoded into the \textit{phase} of the input qubit?
\begin{itemize}
    \item[A)] The entanglement between the two qubits.
    \item[B)] The Phase Kickback effect, which occurs because the ancilla qubit is initialized in the $|-\rangle$ state.
    \item[C)] The final measurement operation collapsing the wavefunction.
\end{itemize}

\vspace{0.3cm}

\textbf{Q3. Reading the Output}\\
You run Deutsch's algorithm on a mysterious Oracle and measure the input qubit at the very end. The result is 100\% certainty of reading \textbf{"0"}. According to the interference step, what did you just prove about the hidden function?
\begin{itemize}
    \item[A)] The function is Balanced (e.g., $f(0) \neq f(1)$).
    \item[B)] The function is Constant (e.g., $f(0) = f(1)$).
    \item[C)] The Oracle contains a hardware error.
\end{itemize}

\vspace{1cm}
\hrule
\vspace{0.2cm}
\footnotesize{\textbf{Solutions:} Q1: C (As stated in QC12); Q2: B (The minus sign is required to kick back the phase); Q3: B (Constructive interference on $|0\rangle$).}

\end{document}